\section{Introduction}
\label{sec:introduction}

World models learn internal states and their dynamics to anticipate how an observed environment may evolve. In visual learning, this approach connects perception with prediction and planning~\citep{assran2025vjepa2,zhou2025dinowm}. Medicine presents a compelling setting for such models: repeated examinations record a patient whose findings may persist, resolve, or progress. A medical world model should encode the evidence available at one examination and predict the state that a later examination may reveal. This requires learning both what constitutes a medically informative state and how that state evolves across observations.

The choice of state is central to this problem. Medical images provide spatial evidence, while reports express findings and their clinical interpretation. Follow-up studies capture changes in both, alongside variation in acquisition and description. Reproducing a future image or report does not by itself specify which aspects of the patient the dynamics model should represent. A useful predictive space must connect clinical meaning with visual evidence and make observations from different time points comparable. The challenge is to construct a latent medical world whose states have clinical and spatial content, and whose transitions can be learned from real patient follow-ups.

Joint-embedding predictive architectures (JEPAs) offer a learning principle for this world. By predicting target embeddings, I-JEPA and V-JEPA learn structure directly in representation space~\citep{assran2023ijepa,bardes2024vjepa}. VLA-JEPA further couples VLM-derived tokens with a latent world model for future visual-state prediction~\citep{sun2026vlajepa}. Medical world modeling has also begun to address patient evolution: MeWM uses generative tumor simulation, CLARITY learns treatment-conditioned latent trajectories, and Clin-JEPA combines EHR dynamics with reusable risk representations~\citep{yang2025medicalworldmodel,ding2025clarity,yang2026clinjepa}. These advances motivate a complementary direction: a JEPA-based medical world model whose predictive state is grounded jointly in images and clinical language, with clinical and spatial tasks shaping the content on which its dynamics operate.

We introduce \method{}, which couples a multimodal state encoder with horizon-conditioned dynamics in a shared latent space. The current image and available report are encoded into state tokens; the world model predicts their evolution toward the state encoded from an observed follow-up image--report pair. This gives JEPA a medical prediction target that combines visual evidence with clinical semantics. To construct that target space, learned readouts aggregate multiple depths of vision--language fusion and the VLM's native vision encoder into four fusion slots and four visual slots. State construction precedes both the task request and the forecast horizon, so the same medical observation defines the starting point for different readouts and temporal predictions.

Learning this world requires supervision for both state content and state evolution. We first train the state encoder and decoders with clinical and spatial tasks, giving the latent space content that can be read out for medical image understanding. Longitudinal training then teaches a predictor to map the current state and a requested horizon to the encoded follow-up. The prediction loss also updates the current-state encoding path, allowing the demands of temporal modeling to shape the states themselves. Future images and reports supply detached state targets and supervised forecast losses; they are unavailable as inference evidence. The two stages thus connect task-supervised state construction with JEPA-style learning of observational patient dynamics.

This world model provides two forms of inference. Its predicted future state augments a native VLM decoder that retains the current image and available clinical history when predicting follow-up findings and reports. Its current state supports classification, VQA, report generation, phrase grounding, segmentation, and super-resolution. For the dense tasks, visual slots condition decoders that also receive the input image, preserving direct access to spatial detail. We assess forecasting across clinical status, disease progression, report fidelity, and probabilistic reliability (\cref{tab:future_prediction}), and current-state reuse through the six downstream tasks (\cref{tab:downstream_tasks}). The downstream protocol withholds the same-exam report and plans matched task supervision and decoder budgets for a no-slots VLM baseline. The tables report available public-model baselines and forecast-pilot results; evaluation of the complete proposed architecture remains pending.

Our contributions are:
\begin{itemize}
    \item A JEPA-based medical world model that learns horizon-conditioned patient-state dynamics by predicting the joint representation of real follow-up images and reports.
    \item A multi-depth multimodal world state with complementary vision--language and visual slots, shaped by clinical and spatial task supervision.
    \item A two-stage learning procedure that establishes the world state through task readout and then adapts its encoding through longitudinal prediction, connecting future forecasting with current-state reuse.
\end{itemize}
